% !TEX encoding = UTF-8
% !TEX program = lualatex
\printnomenclature[3.5cm] % Значение ширины столбца с обозначениями стоит подбирать вручную
% \chapter*{СПИСОК СОКРАЩЕНИЙ И УСЛОВНЫХ ОБОЗНАЧЕНИЙ}
% \addcontentsline{toc}{chapter}{СПИСОК СОКРАЩЕНИЙ И УСЛОВНЫХ ОБОЗНАЧЕНИЙ}
% \markboth{СПИСОК СОКРАЩЕНИЙ И УСЛОВНЫХ ОБОЗНАЧЕНИЙ}{СПИСОК СОКРАЩЕНИЙ И УСЛОВНЫХ ОБОЗНАЧЕНИЙ}

\begin{acronym}
\acro{APT}{Advanced Persistent Threat — целевая многоэтапная атака}
\acro{EDR}{Endpoint Detection and Response — обнаружение и реагирование на конечных устройствах}
\acro{SIEM}{Security Information and Event Management — управление информацией и событиями безопасности}
\acro{SOAR}{Security Orchestration, Automation and Response — оркестрация, автоматизация и реагирование в области безопасности}
\acro{TTP}{Tactics, Techniques, Procedures — тактики, техники и процедуры}
\acro{C2}{Command and Control — команда и управление}
\acro{IDS}{Intrusion Detection System — система обнаружения вторжений}
\acro{IPS}{Intrusion Prevention System — система предотвращения вторжений}
\acro{IOC}{Indicator of Compromise — индикатор компрометации}
\acro{STIX}{Structured Threat Information eXpression — структурированное выражение информации об угрозах}
\acro{TAXII}{Trusted Automated eXchange of Indicator Information — доверенный автоматизированный обмен информацией об индикаторах}
\acro{UCOS}{Unified Cyber Observation Schema — унифицированная схема кибернаблюдений}
\acro{MITRE}{Military Standard — организация, разрабатывающая стандарты безопасности}
\acro{D3FEND}{MITRE D3FEND — фреймворк защиты от кибератак}
\acro{CALDERA}{Cyber Adversary Language and Decision Engine for Red Team Automation — платформа эмуляции атак}
\acro{GNN}{Graph Neural Network — графовая нейронная сеть}
\acro{ML}{Machine Learning — машинное обучение}
\acro{AI}{Artificial Intelligence — искусственный интеллект}
\acro{OS}{Operating System — операционная система}
\acro{AD}{Active Directory — служба каталогов Microsoft}
\acro{DNS}{Domain Name System — система доменных имён}
\acro{HTTP}{Hypertext Transfer Protocol — протокол передачи гипертекста}
\acro{HTTPS}{HTTP Secure — защищённый протокол передачи гипертекста}
\acro{IP}{Internet Protocol — межсетевой протокол}
\acro{URL}{Uniform Resource Locator — унифицированный указатель ресурса}
\acro{PDF}{Portable Document Format — переносимый формат документов}
\acro{VM}{Virtual Machine — виртуальная машина}
\acro{SOC}{Security Operations Center — центр мониторинга безопасности}
\acro{FPR}{False Positive Rate — уровень ложных срабатываний}
\acro{TTD}{Time to Detect — время до обнаружения}
\acro{UC}{Use Case — сценарий использования}
\end{acronym}